\documentclass[11pt,a4paper]{article}
\usepackage[T1]{fontenc}
\usepackage[utf8]{inputenc}
\usepackage{lmodern}
\usepackage{amsmath,amssymb,amsthm,mathtools}
\usepackage[margin=26mm,headheight=15pt]{geometry}
\usepackage{microtype}
\usepackage{booktabs,tabularx,array,longtable}
\usepackage{enumitem}
\usepackage{xcolor}
\usepackage{listings}
\usepackage{fancyhdr}
\usepackage[hyphens]{url}
\usepackage[colorlinks=true,linkcolor=blue!45!black,citecolor=blue!45!black,urlcolor=blue!45!black]{hyperref}
\hypersetup{pdftitle={Cubic Obstructions to Zero-Sector Reduction and Arithmetic Defect Formulas},pdfauthor={Research draft prepared for Keio},pdfsubject={Exact proofs, counterexamples, and explicit priority audit}}
\definecolor{ink}{HTML}{213C50}
\definecolor{muted}{HTML}{52616B}
\definecolor{rulegray}{HTML}{C6CDD2}
\pagestyle{fancy}
\fancyhf{}
\fancyhead[L]{\small\sffamily\color{muted}Cubic obstructions and arithmetic defects}
\fancyhead[R]{\small\sffamily\color{muted}Research note}
\fancyfoot[C]{\small\thepage}
\renewcommand{\headrulewidth}{0.3pt}
\setlength{\emergencystretch}{2em}
\setlength{\parskip}{3pt}
\setlist{itemsep=3pt,topsep=5pt}
\widowpenalty=10000
\clubpenalty=10000
\raggedbottom
\newtheorem{theorem}{Theorem}[section]
\newtheorem{lemma}[theorem]{Lemma}
\newtheorem{proposition}[theorem]{Proposition}
\newtheorem{corollary}[theorem]{Corollary}
\theoremstyle{definition}
\newtheorem{definition}[theorem]{Definition}
\theoremstyle{remark}
\newtheorem{remark}[theorem]{Remark}
\DeclareMathOperator{\Disc}{Disc}
\DeclareMathOperator{\rad}{rad}
\newcommand{\N}{\mathbb N}
\newcommand{\R}{\mathbb R}
\newcommand{\C}{\mathbb C}
\newcommand{\Z}{\mathbb Z}
\newcolumntype{L}[1]{>{\raggedright\arraybackslash}p{#1}}
\newcolumntype{Y}{>{\raggedright\arraybackslash}X}
\lstset{language=Python,basicstyle=\ttfamily\footnotesize,breaklines=true,columns=fullflexible,keepspaces=true,showstringspaces=false,frame=single,rulecolor=\color{rulegray},framesep=6pt,aboveskip=10pt,belowskip=10pt}
\title{\LARGE\bfseries\color{ink}Cubic Obstructions to Zero-Sector Reduction\\[5pt]
\Large and Arithmetic Defect Formulas}
\author{\normalsize Research draft prepared for Keio}
\date{\normalsize 7 November 2026 (UTC)}
\begin{document}
\maketitle
\begin{abstract}
We give an exact negative answer to Question~6.1 of Cardon, Forg\'acs,
Piotrowski, Sorensen and White (2018).  For
$T_{\alpha,\beta}(z^k)=e^{-\alpha k^\beta}z^k$, the original proposed
sufficiency of $\beta\ge2$ fails already at $\beta=3$ and
$\alpha=\log(100/99)$: the image of $(z-1)^3$ has a strictly negative
cubic discriminant.  A general calculation gives an explicit
third-finite-difference obstruction for diagonal exponential semigroups,
and consequently a quadratic-generator rigidity theorem.  Separately,
an odd $p$-adic logarithm expansion proves the Wolstenholme criterion
printed in Me\v{s}trovi\'c's 2015 Remark~1.1, including all small primes.
We identify why this arithmetic conclusion is already an elementary
consequence of earlier established congruences.  Appendices prove the
unchanged binomial power-sum assertion in his 2014 Conjecture~19, derive
an exact local valuation refinement, and give a counterexample to the
first assertion of Conjecture~11.  All mathematical proofs are
unconditional.  Historical priority is not certified, and no result is
claimed to satisfy all six novelty gates specified in the research
brief.  No proof of RH, BSD, abc, or infinitude of Wieferich or
Wolstenholme primes is obtained.
\end{abstract}

\begin{center}
\begin{minipage}{0.92\textwidth}
\small\textbf{Reading this draft.}
``Answer to a printed assertion'' describes a proved mathematical
conclusion.  It does not, by itself, assert that the assertion was still
open on the date of this draft or that the proof is historically new.
Section~\ref{sec:audit} records these distinctions explicitly.  The
single \texttt{.tex} source contains all proofs, references, and an
exact-arithmetic verification program; no external figures, bibliography
files, or source fragments are required to compile it.
\end{minipage}
\end{center}
\vspace{4pt}
\noindent\textbf{Keywords.} Zero-sector reduction; multiplier sequences;
cubic discriminant; finite differences; Wolstenholme primes;
harmonic congruences; power sums.

\clearpage
\tableofcontents
\clearpage

\section{Targets, claims, and conventions}
\label{sec:targets}

The principal target is the parameter criterion proposed in
\cite[Question~6.1]{CardonEtAl2018}.  Our counterexample uses the original
diagonal operator and an admissible real polynomial, without replacing
the target by a newly defined invariant.  The quantitative semigroup
problem studied afterward is auxiliary: its solution explains the
counterexample but is not counted as an additional pre-existing named
problem.

The arithmetic targets are identified by their original numbered
locations.  The following inventory distinguishes the conclusions from
their research-status assessment.

\begin{center}
\small
\begin{tabularx}{\textwidth}{@{}L{0.30\textwidth}Y@{}}
\toprule
Original location & Exact conclusion in this draft\\
\midrule
Cardon et al. (2018), Question 6.1 & Negative answer to the proposed
``$\beta\ge2$'' sufficiency, with a rational cubic certificate.\\[4pt]
Me\v{s}trovi\'c (2015), Remark 1.1 & The displayed harmonic congruence is
equivalent to the Wolstenholme condition.  This follows from earlier
arithmetic and is not claimed as a new method.\\[4pt]
Me\v{s}trovi\'c (2014), Conjecture 19 & The stated congruence holds for
every required $k$ and even $n$; an exact local valuation is also proved.\\[4pt]
Me\v{s}trovi\'c (2014), Conjecture 11 & The first nonexistence assertion
is false at $n=5$.  The second assertion is not resolved here.\\
\bottomrule
\end{tabularx}
\end{center}

Throughout, $v_p$ denotes the usual valuation on $\mathbb Q$, with
$v_p(0)=+\infty$.  A rational congruence $u\equiv v\pmod{p^r}$ means
$v_p(u-v)\ge r$.  When denominators are units at $p$, this is the usual
congruence in $\mathbb Z_{(p)}$.  We write
$\Delta g(m)=g(m+1)-g(m)$.  Zeros are counted with multiplicity;
repeated zeros are allowed in all polynomial tests below.
\section{A negative answer to the power-exponent question of Cardon et al.}
\label{sec:cardon-counterexample}

\subsection{The original problem and its quantifiers}

Following Cardon, Forg\'acs, Piotrowski, Sorensen and White
\cite{CardonEtAl2018}, set
\[
 S(\theta)=\{0\}\cup
 \{z\in\mathbb C\setminus\{0\}:|\arg z|\leq\theta\},
 \qquad 0\leq\theta<\frac{\pi}{2}.
\]
These are sectors about the nonnegative real axis.  Writing
$Z(f)$ for the multiset of zeros of a nonzero polynomial $f$, the
sector-reduction property is
\begin{equation}\label{eq:sector-quantifiers}
 \begin{split}
 &\text{for every }\theta\in(0,\pi/2)
 \text{ there exists }\gamma\in[0,\theta)\text{ such that}\\
 &\qquad Z(f)\subseteq S(\theta)
 \quad\Longrightarrow\quad Z(Tf)\subseteq S(\gamma)
 \quad\text{for every nonzero }f\in\mathbb R[z].
 \end{split}
\end{equation}
All operators considered below have strictly positive diagonal entries,
so $Tf$ is nonzero whenever $f$ is nonzero.

Question~6.1 of \cite{CardonEtAl2018} asks whether, for a positive
parameter $\alpha$, the operator (writing $\beta$ for the exponent
called $p$ in the original question)
\begin{equation}\label{eq:cardon-original-operator}
 T_{\alpha,\beta}[z^k]=\exp(-\alpha k^\beta)z^k,
 \qquad k=0,1,2,\ldots,
\end{equation}
is zero-sector reducing if and only if $\beta\geq2$.
The question appears on page~489 of the published paper and page~10
of arXiv:1802.02641v1.  In particular, its proposed sufficient direction
would make \eqref{eq:sector-quantifiers} hold for every
$\alpha>0$ and every $\beta\geq2$.  The next theorem disproves this
direction with the original operator and the original sector convention.

\begin{theorem}[An exact cubic counterexample]
\label{thm:cardon-negative}
For every
\[
 0<\alpha\leq\log\frac{100}{99},
\]
the polynomial $T_{\alpha,3}[(z-1)^3]$ has two nonreal zeros.
Consequently $T_{\alpha,3}$ is not zero-sector reducing, and
Question~6.1 of \cite{CardonEtAl2018} has a negative answer.
\end{theorem}

\begin{proof}
Put $q=e^{-\alpha}$, so $99/100\leq q<1$.  The image of the
polynomial with a triple positive root is
\[
 P_q(z)=q^{27}z^3-3q^8z^2+3qz-1.
\]
For $az^3+bz^2+cz+d$, the discriminant is
\[
 b^2c^2-4ac^3-4b^3d-27a^2d^2+18abcd.
\]
Substitution and factorization give the exact identity
\begin{equation}\label{eq:cardon-exact-disc}
 \operatorname{Disc}(P_q)
 =-27q^{18}(1-q^6)^2
   \bigl(q^{24}+2q^{18}+3q^{12}-2q^6-3\bigr).
\end{equation}
For completeness, putting $u=q^6$ reduces the factorization to
\[
 3-4u-4u^2+6u^3-u^6
 =-(u-1)^2(u^4+2u^3+3u^2-2u-3).
\]
Bernoulli's inequality gives $q^r\geq1-r/100$ for each
positive integer $r$, while $q^6\leq1$.  Therefore
\begin{align*}
 q^{24}+2q^{18}+3q^{12}-2q^6-3
 &\geq
 \left(1-\frac{24}{100}\right)
 +2\left(1-\frac{18}{100}\right)
 +3\left(1-\frac{12}{100}\right)-2-3\\
 &=\frac1{25}>0.
\end{align*}
Since $0<q<1$, equation~\eqref{eq:cardon-exact-disc} is strictly
negative.  A real cubic with negative discriminant has one real
zero and one nonreal conjugate pair.

Let $\zeta$ be one of the nonreal zeros.  The input $(z-1)^3$
has all its zeros in $S(\theta)$ for every $\theta>0$.
Choose
\[
 0<\theta<\min\{|\arg\zeta|,\pi/2\}.
\]
Then $\zeta\notin S(\theta)$, and hence it cannot belong to any
$S(\gamma)$ with $0\leq\gamma<\theta$.  This contradicts
\eqref{eq:sector-quantifiers}.  The argument uses positive
half-angles and is independent of the convention at $\theta=0$.
\end{proof}

\begin{remark}[A rational certificate]
At the particular value $\alpha=\log(100/99)$, every coefficient
of $P_q$ is rational.  The proof consists of a rational polynomial
identity and an explicit rational strict inequality.  No numerical
root approximation or finite-height exclusion enters the argument.
\end{remark}

\subsection{A quantitative third-difference obstruction}

Let $g:\mathbb N_0\to\mathbb R$ be any real sequence, and set
\begin{equation}\label{eq:diagonal-flow}
 T_t^g[z^k]=e^{-t g(k)}z^k,\qquad t\ge0.
\end{equation}
For a fixed $m\ge0$, normalize its four relevant values by
\[
 c_m=\min_{0\le i\le3}g(m+i),\qquad
 h_i=g(m+i)-c_m,\qquad
 M=\max_{0\le i\le3}h_i.
\]
Thus $0\le h_i\le M$.  Subtracting $c_m$ multiplies the cubic
factor of the transformed polynomial by a positive scalar and does
not change its zeros.  Its third difference is unchanged:
\[
 d=\Delta^3g(m)=h_3-3h_2+3h_1-h_0.
\]

\begin{lemma}[Quantitative cubic obstruction]
\label{lem:quantitative-cubic}
If $d\ne0$, then $M>0$, and
$T_t^g[z^m(z-1)^3]$ has two nonreal zeros for every
\begin{equation}\label{eq:explicit-time}
 0<t<\frac{d^2}{192M^3}.
\end{equation}
More precisely, the normalized cubic
\[
 P_t(z)=e^{-t h_3}z^3-3e^{-t h_2}z^2+3e^{-t h_1}z-e^{-t h_0}
\]
satisfies, for every $t\ge0$,
\begin{equation}\label{eq:quantitative-disc}
 \Disc(P_t)\le -27d^2t^2+5184M^3t^3.
\end{equation}
\end{lemma}

\begin{proof}
The exact discriminant is
\begin{align}
 D(t)={}&81e^{-t(2h_1+2h_2)}
 -108e^{-t(3h_1+h_3)}
 -108e^{-t(h_0+3h_2)}\notag\\
 &-27e^{-t(2h_0+2h_3)}
 +162e^{-t(h_0+h_1+h_2+h_3)}.\label{eq:five-weights}
\end{align}
Expanding through second order gives the polynomial identities
\[
 D(0)=D'(0)=0,\qquad D''(0)=-54d^2.
\]
For instance, the second identity of nonzero degree is equivalent to
\begin{align*}
 &81(2h_1+2h_2)^2-108(3h_1+h_3)^2-108(h_0+3h_2)^2\\
 &\qquad -27(2h_0+2h_3)^2
 +162(h_0+h_1+h_2+h_3)^2=-54d^2.
\end{align*}
Every exponential rate in \eqref{eq:five-weights} belongs to $[0,4M]$,
and the sum of the absolute numerical weights is
$81+108+108+27+162=486$.  It follows that
\[
 |D^{(3)}(s)|\le486(4M)^3=31104M^3\qquad(s\ge0).
\]
Taylor's theorem with its third-order remainder therefore gives
\[
 \bigl|D(t)+27d^2t^2\bigr|\le5184M^3t^3.
\]
This proves \eqref{eq:quantitative-disc}; the right side is strictly
negative under \eqref{eq:explicit-time}.  Finally,
\[
 T_t^g[z^m(z-1)^3]=e^{-t c_m}z^mP_t(z),
\]
so the nonreal cubic zeros are zeros of the original transformed
polynomial.  Every constant in the bound is explicit.
\end{proof}

\begin{corollary}[Every superquadratic power fails near zero]
\label{cor:superquadratic-power}
For every real $\beta>2$, let
\[
 d_\beta=3^\beta-3\cdot2^\beta+3.
\]
Then $d_\beta>0$, and $T_{\alpha,\beta}[(z-1)^3]$ has two nonreal
zeros whenever
\[
 0<\alpha<\frac{d_\beta^2}{192\,3^{3\beta}}.
\]
\end{corollary}
\begin{proof}
The repeated finite-difference formula gives
\[
 d_\beta=\beta(\beta-1)(\beta-2)
 \int_{[0,1]^3}(u+v+w)^{\beta-3}\,du\,dv\,dw>0.
\]
For $2<\beta<3$, first apply the formula to $(\varepsilon+x)^\beta$
and let $\varepsilon\downarrow0$.  The integrand is locally integrable
at the origin; monotone convergence proves the identity.  For
$\beta\ge3$, the formula applies directly.  The lemma with $m=0$,
$g(k)=k^\beta$, and $M=3^\beta$ now proves the assertion.
\end{proof}

\begin{proposition}[The size of the emerging imaginary parts]
\label{prop:splitting}
Under the hypotheses of Lemma~\ref{lem:quantitative-cubic}, the two
nonreal zeros have imaginary parts of magnitude
\[
 \frac{\sqrt3}{2}|d|^{1/3}t^{1/3}+O(t^{2/3})
 \qquad(t\downarrow0).
\]
Their arguments have the same leading magnitude.
\end{proposition}
\begin{proof}
Put $t=u^3$ and $z=1+uw$.  The expression
\[
 Q_u(w)=u^{-3}P_{u^3}(1+uw)
\]
extends analytically to $u=0$, because $P_0(z)=(z-1)^3$.
Its value there is $Q_0(w)=w^3-d$.  Since $d\ne0$, these three
roots are simple, so the analytic implicit-function theorem gives
$w_j(u)=\omega_j+O(u)$, where $\omega_j^3=d$.  Thus
$z_j(t)=1+t^{1/3}\omega_j+O(t^{2/3})$.  The two nonreal cube roots
of the nonzero real number $d$ have imaginary magnitude
$(\sqrt3/2)|d|^{1/3}$.  Finally,
$\arg(1+\eta)=\operatorname{Im}\eta+O(|\eta|^2)$ near $\eta=0$.
\end{proof}

The exact rational example above is stronger than the conservative
general time bound in this particular case.  Neither result asserts
failure for every positive time when $\beta>2$.  The assertion is
failure for all sufficiently small positive times, plus the separately
certified value $\alpha=\log(100/99)$ at $\beta=3$.

\subsection{Rigidity of diagonal semigroups}

\begin{theorem}[Quadratic rigidity]
\label{thm:diagonal-rigidity}
Let $g:\mathbb N_0\to\mathbb R$.  The following are equivalent.
\begin{enumerate}
\item For every $t\geq0$, $T_t^g$ preserves polynomials whose
zeros all lie in $[0,\infty)$.
\item There exists a sequence $t_j>0$, $t_j\to0$, such that
each $T_{t_j}^g$ preserves polynomials whose zeros all lie in
$[0,\infty)$.
\item There are real constants $A,B,C$, with $A\geq0$, such that
\[
 g(k)=Ak^2+Bk+C\qquad(k\in\mathbb N_0).
\]
\end{enumerate}
Furthermore, $T_t^g$ is zero-sector reducing for every $t>0$
if and only if $g(k)=Ak^2+Bk+C$ with $A>0$.
\end{theorem}

\begin{proof}
The first assertion implies the second.  Suppose the second holds.
If $\Delta^3g(m)\ne0$ for some $m$, then
Lemma~\ref{lem:quantitative-cubic} produces nonreal zeros for
$T_{t_j}^g[z^m(z-1)^3]$ whenever $j$ is large enough.
This is impossible.  Thus $\Delta^3g(m)=0$ for every $m\geq0$,
and successive summation of finite differences gives
\[
 g(k)=g(0)+k\Delta g(0)+\binom{k}{2}\Delta^2g(0).
\]
This is a quadratic polynomial $Ak^2+Bk+C$ on $\mathbb N_0$.

The image of $(z-1)^2$ has discriminant
\[
 4e^{-2tg(1)}\left(1-e^{-t\Delta^2g(0)}\right)
 =4e^{-2tg(1)}(1-e^{-2At}).
\]
Nonnegativity for even one $t_j>0$ forces $A\geq0$.

Conversely, the multiplier $e^{-At k^2}$ preserves nonnegative
real-rootedness for $A\geq0$; this is the classical Gaussian
multiplier property, also supplied by
\cite[Theorem~2.5]{CardonEtAl2018} at $\theta=0$.
Multiplication by $e^{-tC}$ and substitution $z\mapsto e^{-tB}z$
preserve nonnegative real-rootedness, which proves the first assertion.

If every $T_t^g$, $t>0$, is zero-sector reducing, it also preserves
$S(0)$.  Indeed, apply sector reduction to every arbitrarily small
positive $\theta$ and use $\bigcap_{\theta>0}S(\theta)=S(0)$.
The established equivalence therefore applies.  If $A=0$,
the operator is merely a nonzero scalar and a positive dilation;
the angles of the roots of $z^2-2(\cos\theta)z+1$ remain
$\pm\theta$, so strict sector reduction fails.

If $A>0$, \cite[Theorem~2.5]{CardonEtAl2018} supplies the uniform
half-angle
\[
 \gamma(t,\theta)
 =\arccos\bigl(\min\{1,e^{At}\cos\theta\}\bigr)<\theta,
 \qquad 0<\theta<\pi/2.
\]
The additional positive dilation and scalar do not alter root
arguments.  This proves strict sector reduction for every $t>0$.
\end{proof}


\begin{remark}[The classification is restricted to semigroups]
Theorem~\ref{thm:diagonal-rigidity} classifies the stated exponential
families under the common small-time preservation hypothesis.  It is
not a characterization of all individual zero-sector reducing
operators, and therefore is not an answer to the broader
Question~6.2 of \cite{CardonEtAl2018}.
\end{remark}

\paragraph{Methodological provenance.}
Repeated-root tests and discriminant geometry have a substantial
history in multiplier-sequence theory; see, for example,
Passare--Rojas--Shapiro \cite{PassareRojasShapiro2011}.  We do not claim
that using a discriminant is a new principle.  The actual calculations
supplied here are the rational sign certificate, the explicit square
$-27(\Delta^3g(m))^2$ in the Taylor coefficient, its uniform remainder
bound, and the consequences with the quantifiers displayed above.
Their historical priority remains unverified.
\section{A printed Wolstenholme criterion and its novelty obstruction}

For a prime $p$, put
\[
 C_p=\binom{2p-1}{p-1},\qquad
 R_j=R_j(p)=\sum_{a=1}^{p-1}a^{-j}.
\]
Rational congruences modulo $p^r$ mean that the difference has $p$-adic
valuation at least $r$.  When all denominators are prime to $p$, this is the
usual congruence in $\mathbb Z_{(p)}$.

Romeo Me\v{s}trovi\'c, \emph{Congruences for Wolstenholme primes},
Czechoslovak Mathematical Journal \textbf{65} (2015), 237--253,
Remark~1.1 on p.~240 \cite{Mestrovic2015}, poses the converse question whether a prime satisfying
\begin{equation}\label{eq:printed-mestrovic}
 C_p\equiv 1+2pR_1+\frac23p^3R_3\pmod{p^7}
\end{equation}
must be a Wolstenholme prime.  The forward implication for Wolstenholme
primes is Corollary~1.1 of the same paper.  The conjecture also occurs as
Remark~1 in arXiv:1108.4178v1 (2011).

\paragraph{The exact conclusion.}
For every prime $p$, \eqref{eq:printed-mestrovic} holds if and only if
$p$ is a Wolstenholme prime, where the latter means
$C_p\equiv1\pmod{p^4}$.  In fact, for every prime $p\ge11$,
\begin{equation}\label{eq:square-defect}
 C_p-1-2pR_1-\frac23p^3R_3
 \equiv 2p^2R_1^2\pmod{p^7}.
\end{equation}

\paragraph{Elementary harmonic valuations.}
If $1\le j<p-1$, then $R_j\equiv0\pmod p$: inversion permutes
$\mathbb F_p^\times$, and the sum of $j$th powers over that group is zero.
For odd $j$ with $j+1<p-1$, reflection gives
\[
 R_j=\sum_{a=1}^{p-1}(p-a)^{-j}
 \equiv -R_j-jpR_{j+1}\pmod{p^2}.
\]
Thus $R_j\equiv0\pmod{p^2}$ for these odd $j$.
In particular, for $p\ge11$,
\[
 v_p(R_1),v_p(R_3),v_p(R_5)\ge2.
\]

\paragraph{Proof of the square defect.}
Define
\[
 F_p(z)=\prod_{a=1}^{p-1}(1+z/a).
\]
For odd $p$, $F_p(p)=C_p$ and
$F_p(-p)=\prod_{a=1}^{p-1}(a-p)/a=1$ exactly.
The $p$-adic logarithm, which is convergent for every factor under
consideration, therefore yields
\begin{equation}\label{eq:odd-log}
 \log C_p
 =\log F_p(p)-\log F_p(-p)
 =2\sum_{\substack{j\ge1\\j\text{ odd}}}\frac{p^j}{j}R_j.
\end{equation}
For $p\ge11$, the $j=5$ summand has valuation at least $7$.
For $j\ge7$, $j-v_p(j)\ge7$, so every remaining term belongs to
$p^7\mathbb Z_p$.  Indeed, the only borderline value $j=7$ has
$v_p(7)=0$, and for $j\ge8$ a positive $v_p(j)$ implies $j\ge11$ and
$v_p(j)\le j-7$.  Consequently
\[
 \log C_p\equiv L:=2pR_1+\frac23p^3R_3\pmod{p^7}.
\]
Here $v_p(L)\ge3$.  In the exponential series all terms of degree
$r\ge3$ have valuation at least $3r-v_p(r!)\ge7$ for $p\ge11$.
For example, $v_p(r!)\le(r-1)/(p-1)$ gives
$3r-v_p(r!)\ge3r-(r-1)/10>7$ for $r\ge3$.
Thus
\[
 C_p\equiv1+L+\frac12L^2\pmod{p^7}.
\]
The product of $2pR_1$ and $(2/3)p^3R_3$ has valuation at least $8$,
and the square of the latter has valuation at least $10$.  This proves
\eqref{eq:square-defect}.

Since $v_p(R_1)\ge2$,
\[
 2p^2R_1^2\equiv0\pmod{p^7}
 \quad\Longleftrightarrow\quad v_p(R_1)\ge3.
\]
The same logarithmic expansion reduced modulo $p^4$ gives
$C_p\equiv1+2pR_1\pmod{p^4}$, so this is precisely the Wolstenholme
condition.  This proves the result for $p\ge11$.

For the remaining primes, the exact defects
$\Delta_p=C_p-1-2pR_1-(2/3)p^3R_3$ are
\[
\begin{array}{c|c|c|c}
p& C_p&\Delta_p&v_p(\Delta_p)\\\hline
2&3&-22/3&1\\
3&10&-81/4&4\\
5&126&15625/2592&6\\
7&1716&50706719/36000&6
\end{array}
\]
None satisfies \eqref{eq:printed-mestrovic}; none is Wolstenholme.
The all-prime assertion follows.

\paragraph{Why this is not a G5 success.}
Me\v{s}trovi\'c, \emph{On the mod $p^7$ determination of
$\binom{2p-1}{p-1}$}, Rocky Mountain Journal of Mathematics
\textbf{44} (2014), 633--648,
DOI 10.1216/RMJ-2014-44-2-633, Theorem~1.1 \cite{Mestrovic2014p7}
(arXiv:1108.1174v1), already proves
\begin{equation}\label{eq:known-mestrovic}
 C_p\equiv1-2pR_1-2p^2R_2+2p^2R_1^2\pmod{p^7}
 \qquad(p\ge11).
\end{equation}
Reflection, without a Wolstenholme assumption, gives
\[
 2R_1+pR_2+p^2R_3+p^3R_4\equiv0\pmod{p^6},
 \qquad 2R_3+3pR_4\equiv0\pmod{p^4}.
\]
To justify the first formula, expand $(p-a)^{-1}$ through degree $5$;
the omitted $p^4R_5$ and $p^5R_6$ terms belong to $p^6\mathbb Z_p$.
For the second formula expand $(p-a)^{-3}$; the omitted
$6p^2R_5$ and $10p^3R_6$ terms belong to $p^4\mathbb Z_p$.
Eliminating $R_4$ gives
\[
 2R_1+pR_2+\tfrac13p^2R_3\equiv0\pmod{p^6}.
\]
Substituting this in \eqref{eq:known-mestrovic} immediately gives
\eqref{eq:square-defect}.  The 2014 paper's Remark~1.6 already uses
the same square $2p^2R_1^2$ to prove a neighboring characterization.
The generic reflection relation also occurs in the proof of Corollary~1.1
of the 2015 paper.  Therefore a new proof of the printed criterion
cannot honestly be advertised as a new noncancellation or height estimate.
An explicit later resolution was not located in the targeted search;
that absence is not a proof of priority.

\section{A realizability obstruction to transferring tower congruences}

Fix an integer $k\ge1$. We call an integer sequence $(a_n)_{n\ge1}$
$k$-congruent if
\[
 a_{mp^r}\equiv a_{mp^{r-1}}\pmod{p^{kr}}
 \qquad(p\text{ prime},\ r\ge1,\ m\ge1,\ p\nmid m).
\]
The following characterization is a direct M\"obius-inversion argument,
of the same type as classical Gauss--Dold congruences. It is included as
a logical obstruction, not as a new solution of a named problem.

\begin{proposition}[M\"obius characterization]
A sequence $(a_n)$ is $k$-congruent if and only if there is an integer
sequence $(b_n)$ such that
\[
 a_n=\sum_{d\mid n}d^k b_d\qquad(n\ge1).
\]
In that case $b_n=n^{-k}\sum_{d\mid n}\mu(n/d)a_d$, so the
sequence $(b_n)$ is unique.
\end{proposition}
\begin{proof}
If the divisor representation holds and $p\nmid m$, then
\[
 a_{mp^r}-a_{mp^{r-1}}
 =p^{kr}\sum_{d\mid m}d^k b_{dp^r},
\]
which proves the congruence. Conversely, define
$c_n=\sum_{d\mid n}\mu(n/d)a_d$. For $p^r\Vert n$, write
$n=p^r m$ with $p\nmid m$. Terms with $p$-exponent less than $r-1$
have zero M\"obius weight, and the remaining terms give
\[
 c_n=\sum_{d\mid m}\mu(m/d)
       (a_{dp^r}-a_{dp^{r-1}}).
\]
Every summand is divisible by $p^{kr}$. This holds for every
$p^r\Vert n$, so $n^k\mid c_n$. Set $b_n=c_n/n^k$ and apply
M\"obius inversion.
\end{proof}

\begin{theorem}[Arbitrary exceptional-prime sets]
Let $E$ be any subset of the prime numbers. For every $k\ge1$ there is
a positive, multiplicative integer sequence $(a_n)$, with $a_1=1$,
such that
\begin{enumerate}
\item $a_n$ is $k$-congruent;
\item for every prime $p$ and every $r\ge1$,
\[
 v_p(a_{p^r}-a_{p^{r-1}})
 =kr+\boldsymbol{1}_{E}(p);
\]
\item for all $n\ge1$, $a_n\le n^{k+2}$.
\end{enumerate}
Consequently, positivity, multiplicativity, polynomial growth, and
uniform tower supercongruences do not, by themselves, imply that the
primes with one additional digit of congruence precision are infinite,
finite, nonempty, or of any prescribed density.
\end{theorem}
\begin{proof}
Put $c_p=p$ if $p\in E$ and $c_p=1$ otherwise. Define a multiplicative
integer sequence $b$ by $b_1=1$ and $b_{p^r}=c_p$ for every $r\ge1$.
Set $a_n=\sum_{d\mid n}d^k b_d$. This is a positive multiplicative
integer sequence, as the Dirichlet convolution of two multiplicative
functions. The proposition proves all the tower congruences, and
\[
 a_{p^r}=1+c_p\sum_{j=1}^r p^{kj},
 \qquad a_{p^r}-a_{p^{r-1}}=c_p p^{kr},
\]
proves the exact valuation statement.
For every $d$, $b_d\le\operatorname{rad}(d)\le d$; hence
\[
 a_n\le\sum_{d\mid n}d^{k+1}
       \le n^{k+1}\tau(n)\le n^{k+2}.
\]
The final logical conclusion follows because $E$ was arbitrary.
\end{proof}

\begin{remark}[Scope]
The sequence constructed in this theorem depends on $E$ and is not
asserted to come from a fixed hypergeometric motive, a fixed Laurent
polynomial, or the actual Fermat or Wolstenholme quotients. The theorem
does not prove independence of any established arithmetic conjecture
from a formal axiom system. Its exact conclusion is that an argument
using only the four displayed properties cannot force a horizontal
prime-distribution conclusion. A successful arithmetic transfer must
use additional structure of its particular sequence.
\end{remark}

\section{What these results do and do not transfer to}
\label{sec:transfer-scope}

\subsection{Zero flows and the Riemann hypothesis}

The real-zero question has a genuine conceptual connection to heat
deformations.  In the standard de Bruijn--Newman framework, the
Riemann hypothesis is equivalent to $\Lambda\le0$.
Rodgers and Tao proved $\Lambda\ge0$ \cite{RodgersTao2020}.
The obstruction in Section~\ref{sec:cardon-counterexample} concerns
arbitrary diagonal semigroups acting on real polynomials.  It provides
no upper bound on $\Lambda$, no uniform estimate for the specific
theta kernel defining the Riemann $\Xi$-function, and no proof that
its zeros are real.  Its useful implication for research design is
precise: replacing a quadratic multiplier exponent by a
superquadratic exponent cannot be justified as a universal
small-time improvement in sector reduction.

\subsection{Exceptional primes and the direction of the abc implication}

For base $2$, a Wieferich prime is an odd prime satisfying
$2^{p-1}\equiv1\pmod{p^2}$.  A Wolstenholme prime satisfies
$C_p\equiv1\pmod{p^4}$.  The criterion proved above is a statement
about each prime individually.  It does not prove that primes
satisfying either exceptional condition occur infinitely often.
The construction in the preceding section isolates the logical gap
between precision at all powers of a fixed prime and distribution
across different primes.

Silverman's theorem obtains infinitude of \emph{non-Wieferich}
primes from the abc conjecture \cite{Silverman1988}.  The negation
in that conclusion matters: it is not a theorem asserting infinitely
many Wieferich primes.  No argument here proves abc, supplies the
required height/radical bound, or removes the hypothesis of that
conditional implication.  Similarly, proving the square-defect
criterion does not produce any new prime with $v_p(R_1)\ge3$.

\subsection{BSD}

The Birch and Swinnerton-Dyer conjecture concerns the agreement of
the arithmetic rank of an elliptic curve with the order of vanishing
of its $L$-function at $s=1$, together with a leading-coefficient
formula in its full form; see the official problem description
\cite{WilesBSD}.  No rank comparison, leading-coefficient formula,
or finiteness theorem for a Tate--Shafarevich group follows from the
polynomial and harmonic calculations in this draft.  A transfer to
BSD would require such additional arithmetic input, which is absent.

\section{Corpus provenance and the six-gate audit}
\label{sec:audit}

\subsection{What was inspected}

The source inventory was pinned to the repository
\emph{Verifiable Mathematical Reasoning Corpus} \cite{VMRC}, commit
\begin{center}
\small\texttt{c30d40119134a59eda6fa5cc6331f27c8e67d034}.
\end{center}
The commit was dated 6 November 2026, 16:24:52 UTC.  The live
inventory contained 130 \texttt{.tex} files.  Of these, 128 matched
an existing local snapshot by Git blob SHA-1, and the two additional
files, numbered 129 and 130, were retrieved separately.  This checks
the identity of the working corpus; it is not a verification of all
its mathematical claims.

The reading was targeted at mechanisms relevant to the requested
problems, including the following documents and themes.
\begin{center}
\small
\begin{tabularx}{\textwidth}{@{}L{0.17\textwidth}L{0.38\textwidth}Y@{}}
\toprule
Corpus item & Theme examined & Role in this draft\\
\midrule
013 & Profinite recognition, Mahler equations, Witt-type escape &
Distinguish finite certificates from a theorem controlling every
remaining obstruction.\\[4pt]
043 & Local certificates, stability, and concavity &
Compare the local valuation method and check that the power-sum
target in Appendix~\ref{app:powers} is not the different
block-reduction assertion discussed there.\\[4pt]
053 & Restriction classes and Baker's $L(1,\chi)$ question &
Identify the need for a genuine nonvanishing term, rather than a
formal equivalence alone.\\[4pt]
112 & Period obstructions for weighted binary expansions &
Separate a prime's lifting valuation from an infinitude claim
over primes.\\[4pt]
126 & Real collisions in the Gaussian heat flow &
Motivate an adversarial preservation test at a repeated real zero.\\[4pt]
129--130 & Certified prime lifting; sparse modular squaring &
Use exact finite arithmetic as a certificate without extrapolating
it to an infinite distribution theorem.\\
\bottomrule
\end{tabularx}
\end{center}

No unverified theorem from the corpus is used as a premise of the
proofs above.  The structural synthesis is methodological: retain
an exact defect, prove that its first surviving term cannot cancel,
and distinguish the resulting local assertion from a global prime
or zero-distribution statement.  This is a description of the
reasoning workflow, not a claim to have invented a historically
new general principle.

\subsection{Gate-by-gate judgment}

The research brief requires all six gates simultaneously.  The
following judgment is deliberately stricter than merely checking a
proof for algebraic correctness.

\begin{longtable}{@{}L{0.09\textwidth}L{0.24\textwidth}L{0.61\textwidth}@{}}
\toprule
Gate & Requirement & Judgment\\
\midrule
\endfirsthead
\toprule Gate & Requirement & Judgment\\\midrule
\endhead
G1 & Pre-existing named target & The cited numbered questions and
conjectures predate this draft.  The auxiliary semigroup
classification and local valuation refinements are not counted as
additional named-problem solutions.\\[5pt]
G2 & Original statement unchanged & The Cardon operator, its
parameters, and its sectors are retained.  The harmonic and
power-sum congruences retain the original constants and moduli.
Conjecture 11 is refuted only in its first assertion.\\[5pt]
G3 & Actually still unresolved & The original publications label
the targets as open questions or conjectures.  This does not certify
their status in November 2026.  For the Wolstenholme criterion,
earlier established arithmetic already yields the result by an
elementary substitution.\\[5pt]
G4 & Unconditional complete proof & The displayed counterexamples,
the general cubic obstruction, the stated semigroup classification,
the harmonic equivalence, and the power-sum theorem have proofs
covering all their quantifiers.  No numerical experiment is used
in place of a universal proof.\\[5pt]
G5 & Substantive new mechanism & An actual negative discriminant
bound and exact local valuation are proved.  However, their
historical methodological novelty is not established.  Repeated-root
discriminant tests are classical, and the Wolstenholme conclusion is
a consequence of known congruences.  No G5 pass is asserted.\\[5pt]
G6 & No earlier duplicate resolution & Targeted searches inspected
the original papers, later indexed references, and related
multiplier-sequence literature.  No exhaustive nonduplication
certificate is available.  Search absence is not used to assert
priority.\\
\bottomrule
\end{longtable}

\noindent\textbf{Overall status.}
There are complete mathematical answers to the printed assertions
specified in this draft.  There are \emph{zero results certified here
to satisfy all of G1--G6}.  In particular, this draft must not be
cited as a verified historically first solution of a major open
problem or as a proof of any of the five headline conjectures in
the research brief.

\subsection{Checks performed and their limits}

The discriminant factorization and its general second-order Taylor
coefficient were checked symbolically.  The rational sign certificate
was checked using exact arithmetic.  Additional modular checks
covered every prime $11\le p<1000$ and $p=16843$, including a
Wolstenholme case; the four smaller primes were evaluated as exact
fractions.  The telescoping identity and its valuation refinement
were checked over finite grids.  These checks supplement the
proofs and help detect transcription errors.  They neither establish
historical priority nor turn a statement about finitely many primes
into an infinitude theorem.

\appendix
\section{The unchanged power-sum Conjecture 19}
\label{app:powers}

Write $S_j(n)=\sum_{a=1}^{n-1}a^j$, including $S_0(n)=n-1$.
Me\v{s}trovi\'c, \emph{On a Congruence Modulo $n^3$ Involving Two
Consecutive Sums of Powers}, Journal of Integer Sequences
\textbf{17} (2014), Article 14.8.4, Conjecture~19 \cite{Mestrovic2014sums}, states that
\begin{equation}\label{eq:conj19}
 \sum_{i=0}^{k}\bigl(1+n(k+1-i)\bigr)
 \binom{2k+2}{2i}S_{2i}(n)\equiv-1\pmod{n^3}
\end{equation}
for every $k\ge1$ and every even positive integer $n$.
The preprint version arXiv:1211.4570v1 (2012) labels this
Conjecture~2.17.  The following is a complete elementary proof of the
unchanged assertion, but no historical priority or G5 success is claimed.

Put $m=2k+2\ge4$, and denote the sum on the left by $T_m(n)$.
The binomial formula gives
\[
 \sum_{\substack{0\le j\le m-2\\ j\text{ even}}}
 \binom mj a^j
 =\frac{(a+1)^m+(a-1)^m}{2}-a^m,
\]
and
\[
 \sum_{\substack{0\le j\le m-2\\j\text{ even}}}
 (m-j)\binom mj a^j
 =\frac m2\bigl((a+1)^{m-1}-(a-1)^{m-1}\bigr).
\]
Summation over $a=1,\ldots,n-1$ telescopes, yielding the exact identity
\begin{equation}\label{eq:telescoping-conj19}
 4\bigl(T_m(n)+1\bigr)
 =(m+2)n^m-2(n-1)^m+mn(n-1)^{m-1}-mn+2.
\end{equation}
Expanding the right side, the coefficients of $n^0,n^1,n^2$ vanish.
For $3\le j<m$, the coefficient of $n^j$ is
$(-1)^j(j-2)\binom mj$; the coefficient of $n^m$ is $2m$.
Hence
\begin{equation}\label{eq:telescoping-coeff}
 \frac{4(T_m(n)+1)}{n^3}
 =\sum_{j=3}^{m-1}(-1)^j(j-2)\binom mj n^{j-3}
   +2m n^{m-3}.
\end{equation}
If $m=2r$, then
\[
 \binom m3=\frac{2r(2r-1)(r-1)}3\equiv0\pmod4,
\]
because $r(r-1)$ is even and division by the odd number $3$ cannot
remove a factor of $4$.  Thus the $j=3$ term of
\eqref{eq:telescoping-coeff} is divisible by $4$.
The $j=4$ term, when present, is $2\binom m4 n$, divisible by $4$
for even $n$.  Every $j\ge5$ term is divisible by $4$ because
$n^{j-3}$ is, and the final term is divisible by $4$ since $m$ is even.
The right side of \eqref{eq:telescoping-coeff} is therefore divisible
by $4$.  Consequently $(T_m(n)+1)/n^3$ is an integer, proving
\eqref{eq:conj19} for all the required $k$ and $n$.

This argument closes the displayed assertion as mathematics.  It uses
an exact telescoping identity followed by elementary coefficient
divisibility, precisely the kind of result excluded by the stated strict
G5 requirement if a new noncancellation/height/zero estimate is required.
Targeted web searches for Conjecture~19 and Conjecture~2.17 did not locate
a later explicit proof; this limited negative finding is not evidence
strong enough to assert G3 or G6.

\paragraph{An exact local valuation strengthening.}
The same expansion proves a more precise noncancellation statement.
Let $\ell$ be a prime dividing $n$.  If $\ell$ is odd, or if
$\ell=2$ and $4\mid n$, then
\begin{equation}\label{eq:exact-local-T}
 v_\ell\bigl(T_m(n)+1\bigr)
 =3v_\ell(n)+v_\ell\binom m3-v_\ell(4).
\end{equation}
This statement is proved here as mathematics; its historical novelty has
not been established.

For $3\le j<m$, the ratio of the absolute value of the degree-$j$
coefficient in \eqref{eq:telescoping-coeff} to the degree-$3$ coefficient
is
\[
 \frac{(j-2)\binom mj}{\binom m3}
 =\frac{6}{j(j-1)}\binom{m-3}{j-3}.
\]
Put $a=v_\ell(n)\ge1$.
For every $j\ge4$, the valuation of the $j$th term relative to the
cubic term is at least
\[
 (j-3)a+v_\ell(6)-v_\ell(j(j-1)).
\]
This is positive for odd $\ell$: when $\ell=3$ use
$v_3(j(j-1))\le j-3$, and when $\ell\ge5$ use
$v_\ell(j(j-1))\le j-4$.
For $\ell=2$ and $a\ge2$, it is at least
$2(j-3)+1-(j-2)=j-3\ge1$.
These elementary bounds follow because at most one of the two consecutive
integers $j,j-1$ is divisible by $\ell$, with the cases $j=4,5$
checked directly and the remaining cases bounded by $\log_\ell j$.

The final term in \eqref{eq:telescoping-coeff} has ratio
\[
 \frac{2m}{\binom m3}n^{m-3}
 =\frac{12}{(m-1)(m-2)}n^{m-3}.
\]
For $m=4$ this is $2n$, of positive valuation for every $\ell\mid n$.
For even $m\ge6$, at most one of $m-1,m-2$ is divisible by $\ell$,
and
$v_\ell((m-1)(m-2))\le\lfloor\log_2(m-1)\rfloor\le m-4$.
The final term therefore also has strictly larger valuation than the
cubic term.  The latter is the unique term of least valuation, proving
\eqref{eq:exact-local-T}.

In particular, for $4\mid n$,
\[
 v_2(T_m(n)+1)
 =3v_2(n)+v_2(m)+v_2(m-2)-3.
\]
Thus the bound modulo $n^3$ is sharp in its $2$-part whenever
$m\equiv4$ or $6\pmod8$ and $4\mid n$.

\section{A counterexample to the first part of Conjecture 11}

The same 2014 paper, Conjecture~11 \cite{Mestrovic2014sums} (arXiv:1211.4570v1,
Conjecture~2.9), asserts that each of
\[
 S_n(n)\equiv\frac{n^3}{2}B_{n-1}\pmod{n^5},\qquad
 S_{n-1}(n)\equiv nB_{n-1}\pmod{n^4}
\]
is satisfied by no integer $n\ge2$.
Here $B_4=-1/30$, and at $n=5$ the first asserted nonexistence fails:
\[
 S_5(5)=1+32+243+1024=1300,\qquad
 \frac{5^3}{2}B_4=-\frac{25}{12},\qquad
 1300+\frac{25}{12}=\frac{5^6}{12}.
\]
The denominator $12$ is a unit modulo $5^5$, so this is an exact
counterexample to the first universal claim, with the original
statement and constants unchanged.

In fact the first congruence holds for every prime $p\ge5$.
For $p\ge7$, Faulhaber's formula writes
\[
 S_p(p)=\sum_{j=1}^{p+1}\frac1j\binom p{j-1}p^jB_{p+1-j}.
\]
The $j=2$ term is $p^3B_{p-1}/2$.  The $j=1$ term is zero.  The $j=3$ term is also zero because $B_{p-2}=0$.
For $4\le j\le p-1$ with a nonzero Bernoulli factor, $j$ is even;
$\binom p{j-1}$ is divisible by $p$, while $j$ and the denominator
of $B_{p+1-j}$ are prime to $p$; the latter assertion follows from the
von Staudt--Clausen denominator theorem, since the nonzero even index
lies strictly between $0$ and $p-1$.  These terms are divisible by $p^5$.
The $j=p$ term involves $B_1$ and has valuation $p\ge7$; the final
$j=p+1$ term has valuation $p+1$.  This proves the assertion for
$p\ge7$, and the displayed calculation proves the remaining $p=5$.

This is a correction to a printed statement, not a claim that a
substantial previously open problem has been solved.  The mechanism is
the classical Faulhaber formula; it does not meet the user's G5 novelty
condition.  The current prior-art search is insufficient to assert that
this printed error has never previously been noticed.


\section{Reproducible exact checks}
\label{app:checks}

Compile this source twice with \texttt{pdflatex}; no shell escape,
external bibliography, images, or included source files are needed.
The following optional verification program uses Python~3 and SymPy.
Save the listing as \texttt{verify.py} and run \texttt{python3 verify.py}.
It performs symbolic identity checks and exact integer/rational
diagnostics; it is not needed for the mathematical proofs.

\begin{lstlisting}
from fractions import Fraction as F
from math import comb, isqrt
from pathlib import Path
import json
import sympy as s

q = s.symbols('q')
z = s.symbols('z')
P = q**27*z**3 - 3*q**8*z**2 + 3*q*z - 1
disc = s.discriminant(P, z)
factor = -27*q**18*(1-q**6)**2*(q**24+2*q**18+3*q**12-2*q**6-3)
assert s.expand(disc-factor) == 0
q0 = s.Rational(99, 100)
bracket = q0**24+2*q0**18+3*q0**12-2*q0**6-3
assert bracket >= s.Rational(1, 25)
assert factor.subs(q, q0) < 0

h0, h1, h2, h3 = s.symbols('h0 h1 h2 h3')
weights = [81, -108, -108, -27, 162]
rates = [2*h1+2*h2, 3*h1+h3, h0+3*h2,
         2*h0+2*h3, h0+h1+h2+h3]
delta = h3-3*h2+3*h1-h0
assert sum(weights) == 0
assert s.expand(sum(c*r for c,r in zip(weights,rates))) == 0
assert s.expand(sum(c*r*r for c,r in zip(weights,rates))/2+27*delta**2) == 0
assert sum(map(abs, weights))*4**3//6 == 5184

def prime(n):
    return n >= 2 and all(n % d for d in range(2,isqrt(n)+1))

def vp(n,p):
    if isinstance(n,F):
        return vp(n.numerator,p)-vp(n.denominator,p)
    if n == 0:
        return float('inf')
    v=0
    while n % p == 0:
        n//=p
        v+=1
    return v

prime_cases = [p for p in range(11,1000) if prime(p)]+[16843]
for p in prime_cases:
    mod=p**7
    inv=[pow(k,-1,mod) for k in range(1,p)]
    r1=sum(inv)%mod
    r3=sum(pow(a,3,mod) for a in inv)%mod
    C=comb(2*p-1,p-1)%mod
    residual=(C-1-2*p*r1-2*pow(3,-1,mod)*p**3*r3)%mod
    assert (residual-2*p*p*r1*r1)%mod == 0
    assert (residual == 0) == (C%p**4 == 1)
    if p == 16843:
        assert residual == 0

small_expected={2:F(-22,3),3:F(-81,4),5:F(15625,2592),7:F(50706719,36000)}
small_records=[]
for p,expected in small_expected.items():
    r1=sum((F(1,k) for k in range(1,p)),F(0))
    r3=sum((F(1,k**3) for k in range(1,p)),F(0))
    defect=comb(2*p-1,p-1)-1-2*p*r1-F(2,3)*p**3*r3
    assert defect == expected
    assert vp(defect,p)<7
    assert (comb(2*p-1,p-1)-1)%p**4 != 0
    small_records.append({'p':p,'defect':str(defect),'valuation':vp(defect,p)})

telescoping_cases=0
valuation_cases=0
for m in range(4,34,2):
    for n in range(2,42):
        T=sum((1+n*(m-j)//2)*comb(m,j)*sum(a**j for a in range(1,n))
              for j in range(0,m,2))
        rhs=(m+2)*n**m-2*(n-1)**m+m*n*(n-1)**(m-1)-m*n+2
        assert 4*(T+1) == rhs
        telescoping_cases+=1
        if n%2 == 0:
            assert (T+1)%n**3 == 0
        for ell in range(2,n+1):
            if prime(ell) and n%ell == 0 and (ell%2 or n%4 == 0):
                assert vp(T+1,ell) == 3*vp(n,ell)+vp(comb(m,3),ell)-vp(4,ell)
                valuation_cases+=1

assert F(sum(a**5 for a in range(1,5)))-F(5**3,2)*F(-1,30) == F(5**6,12)
report={
 'symbolic_discriminant_identity':True,
 'symbolic_general_Taylor_coefficients':True,
 'exact_q_99_100_certificate':True,
 'harmonic_prime_cases':len(prime_cases),
 'included_Wolstenholme_prime':16843,
 'small_prime_cases':small_records,
 'exact_telescoping_cases':telescoping_cases,
 'exact_local_valuation_cases':valuation_cases,
 'printed_nonexistence_counterexample':5,
 'scope':'Diagnostics supplement the universal proofs; no finite test establishes historical novelty or prime infinitude.'}
Path('verification_results.json').write_text(json.dumps(report,indent=2)+'\n')
print(json.dumps(report,indent=2))
\end{lstlisting}

The verified run reported 165 harmonic prime cases, 600 telescoping
cases, and 765 local valuation cases, together with the four exact
small-prime defects.  A failed assertion indicates a reproducibility
or transcription problem that should be investigated before using
the corresponding formula.

\begingroup
\raggedright
\begin{thebibliography}{99}
\bibitem{CardonEtAl2018}
D. A. Cardon, T. Forg\'acs, A. Piotrowski, E. Sorensen, and J. C. White,
\emph{On zero-sector reducing operators},
J. Math. Anal. Appl. \textbf{468} (2018), no. 1, 480--490.
\href{https://doi.org/10.1016/j.jmaa.2018.08.025}{doi:10.1016/j.jmaa.2018.08.025}.
Preprint: \url{https://arxiv.org/abs/1802.02641}.

\bibitem{PassareRojasShapiro2011}
M. Passare, J. M. Rojas, and B. Shapiro,
\emph{New multiplier sequences via discriminant amoebae},
Moscow Math. J. \textbf{11} (2011), no. 3, 547--560.
\url{https://www.mathnet.ru/eng/mmj432}.
Preprint: \url{https://arxiv.org/abs/1001.4076}.

\bibitem{Mestrovic2015}
R. Me\v{s}trovi\'c, \emph{Congruences for Wolstenholme primes},
Czechoslovak Math. J. \textbf{65} (2015), no. 1, 237--253.
\url{https://dml.cz/dmlcz/144224}.
Preprint: \url{https://arxiv.org/abs/1108.4178}.

\bibitem{Mestrovic2014p7}
R. Me\v{s}trovi\'c,
\emph{On the mod $p^7$ determination of $\binom{2p-1}{p-1}$},
Rocky Mountain J. Math. \textbf{44} (2014), no. 2, 633--648.
\href{https://doi.org/10.1216/RMJ-2014-44-2-633}{doi:10.1216/RMJ-2014-44-2-633}.
Preprint: \url{https://arxiv.org/abs/1108.1174}.

\bibitem{Mestrovic2014sums}
R. Me\v{s}trovi\'c,
\emph{On a congruence modulo $n^3$ involving two consecutive sums of powers},
J. Integer Sequences \textbf{17} (2014), Article 14.8.4.
\url{https://cs.uwaterloo.ca/journals/JIS/VOL17/Mestrovic/mes4.html}.
Preprint: \url{https://arxiv.org/abs/1211.4570}.

\bibitem{RodgersTao2020}
B. Rodgers and T. Tao,
\emph{The de Bruijn--Newman constant is non-negative},
Forum Math. Pi \textbf{8} (2020), e6.
\href{https://doi.org/10.1017/fmp.2020.6}{doi:10.1017/fmp.2020.6}.
Preprint: \url{https://arxiv.org/abs/1801.05914}.

\bibitem{Silverman1988}
J. H. Silverman,
\emph{Wieferich's criterion and the abc-conjecture},
J. Number Theory \textbf{30} (1988), no. 2, 226--237.
\href{https://doi.org/10.1016/0022-314X(88)90019-4}{doi:10.1016/0022-314X(88)90019-4}.

\bibitem{WilesBSD}
A. Wiles, \emph{The Birch and Swinnerton-Dyer conjecture},
official Millennium Prize problem description, Clay Mathematics Institute.
\url{https://www.claymath.org/wp-content/uploads/2022/02/MPPc.pdf}.
Problem page, accessed 7 November 2026:
\url{https://www.claymath.org/millennium/birch-and-swinnerton-dyer-conjecture/}.

\bibitem{VMRC}
QEEDDD, \emph{Verifiable Mathematical Reasoning Corpus},
repository snapshot at commit
\texttt{c30d40119134a59eda6fa5cc6331f27c8e67d034},
6 November 2026.  Research drafts; no blanket verification of their
claims is assumed in this note.
\url{https://github.com/QEEDDD/Verifiable-Mathematical-Reasoning-Corpus/tree/c30d40119134a59eda6fa5cc6331f27c8e67d034}.
\end{thebibliography}
\endgroup
\end{document}
